\section{Kth Smallest Element in an Array}
\label{sec:heap-kth-smallest}

\problemstatement{Given an array $\textit{nums}$ of $n$ integers and an
integer $k$, find the $k$-th smallest element in the array.}

\begin{patternbox}
The mirror image of Section~\ref{sec:heap-kth-largest}: \emph{``Kth
smallest''} calls for the \emph{opposite}-polarity heap, a
\textbf{max}-heap of size $k$. The root of a size-$k$ max-heap is the
\emph{largest} of the current best-$k$-smallest candidates -- exactly the
one to discard the instant a smaller element arrives.
\end{patternbox}

\subsection*{Understanding the problem carefully}

By the same reasoning as Section~\ref{sec:heap-kth-largest}, applied with
every comparison flipped: maintain a max-heap that never holds more than
$k$ elements. Push each new element, then if the heap exceeds size $k$,
pop the maximum (the current weakest member of the ``smallest $k$
so far'' set). After processing the whole array, the heap holds the $k$
smallest elements, and its root -- the largest of those $k$ -- is the
$k$-th smallest element overall.

\begin{ideabox}[Reuse, with Every Comparison Flipped]
Apply Algorithm~\ref{sec:heap-kth-largest}'s \textsc{KthLargest} procedure
verbatim, replacing the min-heap with a max-heap and ``pop the minimum''
with ``pop the maximum.''
\end{ideabox}

\subsection*{Algorithm}

\begin{enumerate}
  \item Initialize an empty max-heap.
  \item For each element $x$ in $\textit{nums}$:
  \begin{enumerate}
    \item Push $x$ onto the heap.
    \item If the heap's size exceeds $k$: pop the maximum.
  \end{enumerate}
  \item Return the heap's current maximum (its root).
\end{enumerate}

\begin{algorithm}[H]
\caption{Kth Smallest Element}
\begin{algorithmic}[1]
\Procedure{KthSmallest}{\textit{nums}, $k$}
  \State $\textit{heap} \gets$ empty max-heap
  \For{each $x$ in \textit{nums}}
    \State push $x$ onto \textit{heap}
    \If{$|\textit{heap}| > k$}
      \State pop the maximum from \textit{heap}
    \EndIf
  \EndFor
  \State \Return top of \textit{heap}
\EndProcedure
\end{algorithmic}
\end{algorithm}

\subsection*{Dry run}

\dryrun{Take $\textit{nums} = [7,10,4,3,20,15]$, $k = 3$.}

\begin{itemize}
  \item Push $7$: heap $=\{7\}$.
  \item Push $10$: heap $=\{7,10\}$.
  \item Push $4$: heap $=\{4,7,10\}$. Size $3 \le 3$.
  \item Push $3$: heap $=\{3,4,7,10\}$. Size $4>3$: pop max ($10$). Heap
        $=\{3,4,7\}$.
  \item Push $20$: heap $=\{3,4,7,20\}$. Size $4>3$: pop max ($20$). Heap
        $=\{3,4,7\}$.
  \item Push $15$: heap $=\{3,4,7,15\}$. Size $4>3$: pop max ($15$). Heap
        $=\{3,4,7\}$.
\end{itemize}

The final heap holds $\{3,4,7\}$, the three smallest elements; its root
(maximum of these three) is $7$, the $3$rd smallest element overall.
Sorting $[3,4,7,10,15,20]$ confirms index $k-1=2$ holds $7$.

\subsection*{C++ implementation}

\begin{lstlisting}
#include <bits/stdc++.h>
using namespace std;

int kthSmallest(vector<int>& nums, int k) {
    priority_queue<int> maxHeap; // max-heap (default in C++)

    for (int x : nums) {
        maxHeap.push(x);
        if ((int)maxHeap.size() > k) {
            maxHeap.pop();
        }
    }
    return maxHeap.top();
}

int main() {
    vector<int> nums = {7, 10, 4, 3, 20, 15};
    cout << "3rd smallest: " << kthSmallest(nums, 3) << endl;
    return 0;
}
\end{lstlisting}

\begin{complexitybox}
\textbf{Time Complexity.}
\[
  O(n \log k),
\]
identical in structure to Section~\ref{sec:heap-kth-largest}.
\textbf{Space Complexity.} $O(k)$.
\end{complexitybox}

\begin{takeawaybox}
Once Kth Largest is internalized, Kth Smallest costs nothing extra to
derive -- it is the same algorithm with every comparison polarity
reversed. The general rule: a size-$k$ heap of the \emph{opposite}
polarity to what you are looking for always keeps the right boundary
element at the root.
\end{takeawaybox}
